% GENERATED FILE. Edit canonical lessons, then run npm run generate:books.
% canonical-source-hash: fnv1a64:f9f3506c6c48ef2e
% canonical-lessons: UR-C184-bimar-parna, UR-C184-sujna, UR-C184-kanpna, UR-C184-bhigna, UR-C184-safar-karna

\chapter{To fall ill, To swell, To shiver, To get wet, To travel}
\label{ch:ur-to-fall-ill-to-swell-to-shiver-to-get-wet-to-travel}
\hlchaptermodality{184}

\begin{chapteropening}
\textbf{By the end of this chapter:} \emph{I can say to fall ill, to swell, to shiver, to get wet and to travel.}

Complete the last lesson of chapter 184: I can say to fall ill, to swell, to shiver, to get wet and to travel.
\end{chapteropening}

\section[\texorpdfstring{\ur{بیمار} \ur{پڑنا} (bīmār paṛnā)}{بیمار پڑنا (bīmār paṛnā)}]{\ur{بیمار} \ur{پڑنا} (bīmār paṛnā) — to fall ill}

\label{lesson:UR-C184-bimar-parna}



\begin{quote}
Before the new one: say the Urdu for to die, then the Urdu for to increase.
\end{quote}

\subsection*{You'll want to know: \ur{بیمار} \ur{پڑنا}}
\textbf{\ur{بیمار} \ur{پڑنا}} — \emph{bīmār paṛnā} — \textquotedblleft{}to fall ill\textquotedblright{}.

\begin{cousinweb}
\textbf{\ur{بیمار} \ur{پڑنا}} (\emph{bīmār paṛnā}) is \textquotedblleft{}to fall ill\textquotedblright{}. \emph{bīmār}, \textquotedblleft{}ill\textquotedblright{}, is usually traced to Persian \emph{bīmār}, and \emph{paṛnā}, \textquotedblleft{}to fall, to happen\textquotedblright{}, comes from Sanskrit \emph{patati}, \textquotedblleft{}flies; falls\textquotedblright{} (CDIAL 7722).
\end{cousinweb}

\begin{grammarlens}[title={What you've built}]
One more word for this chapter.
\end{grammarlens}

\subsection*{Your turn}
\begin{itemize}
\raggedright
  \item \emph{Say it:} \emph{bīmār paṛnā}
  \item \emph{Say it:} \emph{bīmār paṛnā}, once more
  \item \emph{Say it:} say \emph{bīmār paṛnā}
  \item \emph{Recall:} say the Urdu for to die, then the Urdu for to increase, then say \emph{bīmār paṛnā} again
\end{itemize}

\begin{tcolorbox}[breakable,colback=teal!4,colframe=teal!35!black,title={Before you move on}]
What does \ur{بیمار} \ur{پڑنا} mean? (\textquotedblleft{}To fall ill\textquotedblright{}.) Say it once more.
\end{tcolorbox}
\section[\texorpdfstring{\ur{سوجنا} (sūjnā)}{سوجنا (sūjnā)}]{\ur{سوجنا} (sūjnā) — to swell}

\label{lesson:UR-C184-sujna}



\begin{quote}
Before the new one: say the Urdu for to increase, then the Urdu for to fall ill.
\end{quote}

\subsection*{You'll want to know: \ur{سوجنا}}
\textbf{\ur{سوجنا}} — \emph{sūjnā} — \textquotedblleft{}to swell\textquotedblright{}.

\begin{cousinweb}
\textbf{\ur{سوجنا}} (\emph{sūjnā}) is a tadbhava, a word inherited from Sanskrit and reshaped by sound change. It comes from \emph{śūyate}, \textquotedblleft{}is swollen\textquotedblright{} (CDIAL 12568).
\end{cousinweb}

\begin{grammarlens}[title={What you've built}]
One more word for this chapter.
\end{grammarlens}

\subsection*{Your turn}
\begin{itemize}
\raggedright
  \item \emph{Say it:} \emph{sūjnā}
  \item \emph{Say it:} \emph{sūjnā}, once more
  \item \emph{Say it:} say \emph{sūjnā}
  \item \emph{Recall:} say the Urdu for to increase, then the Urdu for to fall ill, then say \emph{sūjnā} again
\end{itemize}

\begin{tcolorbox}[breakable,colback=teal!4,colframe=teal!35!black,title={Before you move on}]
What does \ur{سوجنا} mean? (\textquotedblleft{}To swell\textquotedblright{}.) Say it once more.
\end{tcolorbox}
\section[\texorpdfstring{\ur{کانپنا} (kāṅpnā)}{کانپنا (kāṅpnā)}]{\ur{کانپنا} (kāṅpnā) — to shiver, to tremble}

\label{lesson:UR-C184-kanpna}



\begin{quote}
Before the new one: say the Urdu for to fall ill, then the Urdu for to swell.
\end{quote}

\subsection*{You'll want to know: \ur{کانپنا}}
\textbf{\ur{کانپنا}} — \emph{kāṅpnā} — \textquotedblleft{}to shiver, to tremble\textquotedblright{}.

\begin{cousinweb}
\textbf{\ur{کانپنا}} (\emph{kāṅpnā}) is a tadbhava, a word inherited from Sanskrit and reshaped by sound change. It comes from \emph{kampate}, \textquotedblleft{}trembles, shivers\textquotedblright{} (CDIAL 2767).
\end{cousinweb}

\begin{grammarlens}[title={What you've built}]
One more word for this chapter.
\end{grammarlens}

\subsection*{Your turn}
\begin{itemize}
\raggedright
  \item \emph{Say it:} \emph{kāṅpnā}
  \item \emph{Say it:} \emph{kāṅpnā}, once more
  \item \emph{Say it:} say \emph{kāṅpnā}
  \item \emph{Recall:} say the Urdu for to fall ill, then the Urdu for to swell, then say \emph{kāṅpnā} again
\end{itemize}

\begin{tcolorbox}[breakable,colback=teal!4,colframe=teal!35!black,title={Before you move on}]
What does \ur{کانپنا} mean? (\textquotedblleft{}To shiver, to tremble\textquotedblright{}.) Say it once more.
\end{tcolorbox}
\section[\texorpdfstring{\ur{بھیگنا} (bhīgnā)}{بھیگنا (bhīgnā)}]{\ur{بھیگنا} (bhīgnā) — to get wet}

\label{lesson:UR-C184-bhigna}



\begin{quote}
Before the new one: say the Urdu for to swell, then the Urdu for to shiver.
\end{quote}

\subsection*{You'll want to know: \ur{بھیگنا}}
\textbf{\ur{بھیگنا}} — \emph{bhīgnā} — \textquotedblleft{}to get wet\textquotedblright{}.

\begin{cousinweb}
\textbf{\ur{بھیگنا}} (\emph{bhīgnā}) is a tadbhava, a word inherited and reshaped by sound change. It comes from \emph{*bhiyagna}, \textquotedblleft{}smeared, wet\textquotedblright{}, a form CDIAL reconstructs from its descendants (CDIAL 9500). CDIAL builds it on Sanskrit \emph{añj-}, \textquotedblleft{}to anoint\textquotedblright{}.
\end{cousinweb}

\begin{grammarlens}[title={What you've built}]
One more word for this chapter.
\end{grammarlens}

\subsection*{Your turn}
\begin{itemize}
\raggedright
  \item \emph{Say it:} \emph{bhīgnā}
  \item \emph{Say it:} \emph{bhīgnā}, once more
  \item \emph{Say it:} say \emph{bhīgnā}
  \item \emph{Recall:} say the Urdu for to swell, then the Urdu for to shiver, then say \emph{bhīgnā} again
\end{itemize}

\begin{tcolorbox}[breakable,colback=teal!4,colframe=teal!35!black,title={Before you move on}]
What does \ur{بھیگنا} mean? (\textquotedblleft{}To get wet\textquotedblright{}.) Say it once more.
\end{tcolorbox}
\section[\texorpdfstring{\ur{سفر} \ur{کرنا} (safar karnā)}{سفر کرنا (safar karnā)}]{\ur{سفر} \ur{کرنا} (safar karnā) — to travel}

\label{lesson:UR-C184-safar-karna}



\begin{quote}
Before the new one: say the Urdu for to shiver, then the Urdu for to get wet.
\end{quote}

\subsection*{You'll want to know: \ur{سفر} \ur{کرنا}}
\textbf{\ur{سفر} \ur{کرنا}} — \emph{safar karnā} — \textquotedblleft{}to travel\textquotedblright{}.

\begin{cousinweb}
\textbf{\ur{سفر} \ur{کرنا}} (\emph{safar karnā}) joins \emph{safar}, \textquotedblleft{}a journey\textquotedblright{}, and \emph{karnā}, \textquotedblleft{}to do\textquotedblright{}. \emph{safar} is Arabic, \textquotedblleft{}a journey\textquotedblright{}, and reached Urdu through Persian (Lane, root \emph{s-f-r}); \emph{karnā} comes from Sanskrit \emph{karoti}, \textquotedblleft{}does\textquotedblright{} (CDIAL 2814).
\end{cousinweb}

\begin{grammarlens}[title={What you've built}]
One more word for this chapter.
\end{grammarlens}

\subsection*{Your turn}
\begin{itemize}
\raggedright
  \item \emph{Say it:} \emph{safar karnā}
  \item \emph{Say it:} \emph{safar karnā}, once more
  \item \emph{Say it:} say \emph{safar karnā}
  \item \emph{Recall:} say the Urdu for to shiver, then the Urdu for to get wet, then say \emph{safar karnā} again
\end{itemize}

\begin{tcolorbox}[breakable,colback=teal!4,colframe=teal!35!black,title={Before you move on}]
What does \ur{سفر} \ur{کرنا} mean? (\textquotedblleft{}To travel\textquotedblright{}.) Say it once more.
\end{tcolorbox}
